
\chapter{ניסויים ותוצאות}

\section{פרוטוקול הערכה}

השתמשנו ב-\textenglish{AgentBench} לצורך הערכת מסגרות התיאום \cite{wu2023autogen}. המשימות כוללות: חיפוש מידע, כתיבת קוד ובדיקתו, ניתוח מסמכים ואינטגרציה עם \textenglish{API}ים חיצוניים. כל מסגרת הורצה \textenglish{100} פעמים על כל קטגוריה.

\section{תוצאות}

\begin{figure}[H]
\centering
\includegraphics[width=0.82\textwidth]{assets/results_graph.png}
\caption{\textenglish{Latency: Sequential vs Parallel Orchestration (Python-generated via matplotlib)}}
\label{fig:graph}
\end{figure}

הגרף מדגים את יתרון הביצוע המקבילי: עם \textenglish{16} סוכנים, \textenglish{Parallel Orchestration} מציג זמן ריצה \textenglish{$\sim$6.2×} קצר יותר מ-\textenglish{Sequential} \cite{yao2022react}. עלות האסימון גדלה ב-\textenglish{$\sim$40\%} בשל תקשורת בין-סוכנים, אך \textenglish{throughput} הכולל גדל ב-\textenglish{5.1×}.

\section{ניתוח כשלים}

הגורמים העיקריים לכישלון (\textenglish{18\% error rate ב-ReAct}) הם: \textenglish{hallucination} של פרמטרי \textenglish{API}, לולאות אינסוף (\textenglish{infinite loops}) ועיוות הקשר בשיחות ארוכות. \textenglish{CrewAI} עם מנגנון \textenglish{Watchdog} ותקציב \textenglish{MAX\_ITER} פוחת ל-\textenglish{6\%} כשל \cite{nakash2024skillsieve}.
